\documentclass[10pt,a4paper]{amsart}
\usepackage[margin=19mm]{geometry}
\usepackage{amsmath,amssymb,mathtools}
\usepackage[scr=rsfs,cal=euler]{mathalfa}
\usepackage{xfrac}
\usepackage[unicode,hidelinks,bookmarks=false]{hyperref}
\newcommand{\ie}{i.e.\ }
\newcommand{\cf}{cf.\ }
\newcommand{\resp}{respectively\ }
\newcommand{\Subsection}{\subsection}
\providecommand{\nobreakdash}{\nobreak\hbox{-}}
\setlength{\emergencystretch}{2em}
\multlinegap0pt
\usepackage{bm}
\usepackage{bbm}
\usepackage{dsfont}
\usepackage{xspace}
\usepackage{cancel}
\usepackage{mathtools}
\usepackage[shortlabels]{enumitem}
\usepackage{amsthm}
\makeatletter
\@ifundefined{theorem}{\newtheorem{theorem}{Theorem}}{}
\@ifundefined{definition}{\newtheorem{definition}{Definition}}{}
\@ifundefined{lemma}{\newtheorem{lemma}[theorem]{Lemma}}{}
\makeatother
\makeatletter
\@ifundefined{claimproof}{\newenvironment{claimproof}[1]{\par\noindent\emph{Proof of the claim:}\space#1}{\leqslantavevmode\unskip\penalty9999 \hbox{}\nobreak\hfill\quad\hbox{$\blacksquare$}}}{}
\@ifundefined{theoremproof}{\newenvironment{theoremproof}[1]{\par\noindent\emph{Proof of theorem}\space#1}{\leqslantavevmode\unskip\penalty9999 \hbox{}\nobreak\hfill\quad\hbox{$\square$}}}{}
\makeatother
\title[Tur\'an-type and tiling problems in oriented graphs]{Tur\'an-type and tiling problems in oriented graphs}
\author{Ming Chen, Wenxu Lu, Yun Wang, Zhiwei Zhang}
\date{Source: arXiv:2603.21971v1 (CC BY 4.0)}
\begin{document}
\maketitle

\noindent\emph{Source and licence.} Tur\'an-type and tiling problems in oriented graphs. Version arXiv:2603.21971v1 (CC BY 4.0), distributed under the Creative Commons Attribution 4.0 International licence, as recorded in the arXiv licence metadata for this exact version and shown on the abstract page https://arxiv.org/abs/2603.21971v1. Full rights notice: licenses/arxiv-2603.21971-CC-BY-4.0.txt.\par\smallskip

\noindent\textit{Editorial scope.} Counted unit: the Lemma whose source label is \texttt{LEM-edgegood} in the article (source0.tex lines 1271--1273, source label \texttt{LEM-edgegood}), delivered together with the article's own complete original proof of it (source0.tex lines 1275--1345, one \texttt{proof} environment of 71 lines). No mathematics is written, repaired or re-derived: statement and proof are the article's own text line for line. Required context retained uncounted: EP2 (lines 491--494); LEM:1bcbalance (lines 530--537); LEM:fexbad (lines 1159--1161); LEM:greedily Dabc (lines 1165--1167); DEF:robust (lines 1258--1265). Every definition, display and label of those lines is retained unchanged. Omitted from this package and not redistributed: 1--490, 495--529, 538--1158, 1162--1164, 1168--1257, 1266--1270, 1274--1274, 1346--2093. Nothing inside a retained excerpt is omitted or altered: every retained body is delivered exactly as the article's lines are written, so no omission receipt of any kind accompanies this package. Citations: the retained excerpts carry no citation command at all, so no bibliography entry of the article is transcribed and no reference is redistributed. Reference adaptation: none was needed and none was made; every reference inside the delivered unit resolves inside it, so no unresolved reference or citation is produced. Printed slips kept literal: no printed word, symbol, hypothesis or number of the counted statement or of its proof was altered. Layout and class: the article's own class is \texttt{amsart} with packages including setspace, fullpage, datetime, indentfirst, xcolor, array, enumitem, amsmath, amsfonts, amssymb, graphicx, mathrsfs, hyperref, geometry; the wrapper is the lane's pinned \texttt{amsart} editorial wrapper at 10pt on A4, which restates the theorem-like environments the retained text needs and adds the article's own macro definitions that the retained text needs (none), with extra wrapper packages (bm, bbm, dsfont, xspace, cancel, mathtools, enumitem). The delivered numbering is the unit's own. Assets: no retained span contains a figure environment, a graphics-inclusion command, a PDF-page inclusion, a table or a TikZ picture; no image file is copied into this package, the package declares no asset dependency and no image file is redistributed with it. Licence: version arXiv:2603.21971v1 of this article is distributed under the Creative Commons Attribution 4.0 International licence, as recorded in the arXiv licence metadata for this exact version and shown on the abstract page https://arxiv.org/abs/2603.21971v1. A full rights notice is delivered at licenses/arxiv-2603.21971-CC-BY-4.0.txt. Characters: every retained excerpt is pure ASCII, so the delivered engine, XeLaTeX with its default Latin Modern text font, reports no missing character.
\begin{enumerate}[label=\upshape \textbf{(EP\arabic{enumi})},ref=\upshape (EP\arabic{enumi}),resume=mylist]
    \setlength{\itemindent}{1.5em}
    \item $d^+(v,V_{i+1}),\,d^-(v,V_{i-1})\geqslant n/6-O(\gamma n)$ for each $v\in V_i$.\label{EP2}
\end{enumerate}

\begin{lemma}\label{LEM:1bcbalance}
Given $b,c\in \mathbb{N}$ with  $1\leqslant b\leqslant c$ and $3\mid(1+b+c)$, let
$0<1/n\ll \gamma\ll 1$, where $n\in(1+b+c)\mathbb{N}$.
Suppose that $G$ is an $n$-vertex semi-regular tournament.
If $G$ is $\gamma$-extremal, then $G$ contains a constant-sized $D_{1,b,c}$-tiling $\mathcal{H}$ such that
$V(G)\setminus V(\mathcal{H})$ admits a $\gamma^{1/3}$-superextremal partition $(V_1,V_2,V_3)$ with
$|V_i|\equiv 0\pmod{1+b+c}$ for each $i\in[3]$.
\end{lemma}

\begin{lemma}\label{LEM:fexbad}
    Let $0<1/n\ll \eta \ll\gamma \ll 1$. Suppose that $G$ is an $n$-vertex oriented graph with $\delta^0(G)\geqslant(1/2-\eta)n$. If $G$ is $\gamma$-extremal, then $V(G)$ has a $\gamma^{1/3}$-superextremal partition. 
\end{lemma} 

\begin{lemma}\label{LEM:greedily Dabc}
    Given $a,b,c\in\mathbb{N}$ with $1\leqslant a\leqslant b\leqslant c$, let $0<1/n \ll \eta \ll \beta \ll \gamma \ll 1$. Suppose that $G$ is an $n$-vertex oriented graph with $\delta^0(G)\geqslant(1/2-\eta)n$. If $V(G)$ has a $\gamma$-extremal partition $(V_1,V_2,V_3)$, then for each $i\in[3]$, $G$ contains at least $\beta n$ disjoint copies of $D_{a,b,c}$ of type-$(i,i+1,i+2)$.
\end{lemma}

\begin{definition}\label{DEF:robust}
Given $b,c\in\mathbb{N}$ with $1\leqslant b\leqslant c$, let $(V_1,V_2,V_3)$ be a partition of $V(G)$, and let $i,j\in[3]$ be distinct. 
A set $X\subseteq V(G)$ is called  \emph{$\gamma$-transfer for $(V_i,V_j)$} if, for every set $A\subseteq V(G)\setminus X$ with $|A|=O( \gamma n)$, there exists a set $S\subseteq V(G)\setminus A$ such that
\begin{itemize}
\item $|S|\leqslant 3(1+b+c)$, and $G[S]$ contains a $D_{1,b,c}$-factor;
\item $|V_i\setminus S| \equiv |V_i|-1 \pmod{1+b+c}$ and $|V_j\setminus S| \equiv |V_j|+1 \pmod{1+b+c}$.
\end{itemize}
\end{definition}

\begin{lemma}\label{LEM-edgegood}
    Given $b,c\in\mathbb{N}$ with $1\leqslant b\leqslant c$, let $0<1/n\ll\gamma\ll 1$. Suppose that $G$ is an $n$-vertex semi-regular tournament. If $(V_1,V_2,V_3)$ is a $\gamma$-superextremal partition of $V(G)$, then the vertex set of every edge in $E(V_{i+1},V_i)$ is $\gamma$-transfer for both $(V_i,V_{i-1})$ and $(V_{i+1},V_{i-1})$.
\end{lemma}

\begin{proof}[\textbf{Proof of Lemma \ref{LEM:1bcbalance}}] 
Given $b,c\in\mathbb{N}$ with $1\leqslant b\leqslant c$, set 
$h=1+b+c$. Let $G$ be a semi-regular tournament on $n\in h\mathbb{N}$ vertices that is $\gamma$-extremal. Since $3\mid h$ and $n\in h\mathbb{N}$, it follows that $3\mid n$. By Lemma \ref{LEM:fexbad}, $V(G)$ has a $\gamma^{1/3}$-superextremal partition, say $(V_1,V_2,V_3)$. Relabeling the parts if necessary, we may assume w.l.o.g. that $V_1$ is the largest part. Moreover, by reversing all edges of $G$ and swapping the labels of $V_2$ and $V_3$ if necessary, we may further assume that $|V_2|\geqslant |V_3|$. Set $|V_1|=n/3+g$, $|V_2|=n/3+f$, and $|V_3|=n/3-g-f$, where $g$ and $f$ are integers since $3\mid n$. It follows from $|V_1|\geqslant |V_2|\geqslant |V_3|$ that $g\geqslant 0$ and $g\geqslant f\geqslant -g-f$.

We denote by $M_{21}$ and $M_{13}$ maximum matchings in $E(V_2,V_1)$ and $E(V_1,V_3)$, respectively. Next we claim that $|M_{21}|\geqslant 2g+f-1$ and $|M_{13}|\geqslant g-f-1$. Considering the out-degrees of vertices in $V_2$, we have 
\begin{align*}
    e(V_2,V_1)&\geqslant \sum_{v\in V_2} d^+(v)-e(V_2)-e(V_2,V_3)\\
    &\geqslant |V_2|\lfloor (n-1)/2\rfloor-|V_2|(|V_2|-1)/2-|V_2||V_3|\\
    &\geqslant |V_2|(g+f/2-1/2).
\end{align*}

Let $H$ be the underlying bipartite graph corresponding to the set of edges from $V_2$ to $V_1$. Clearly, $e(H)\geqslant|V_2|(g+f/2-1/2)$. On the other hand, since $(V_1,V_2,V_3)$ is $\gamma^{1/3}$-superextremal, \ref{EP2} implies that every vertex of $V_2$ has degree at most $n/6 + O(\gamma^{1/3}n)$ in $H$. Since $H$ is bipartite, by K\"{o}nig Theorem, the size of a maximum matching equals the size of a minimum vertex cover in $H$. Thus, we obtain that
$$
e(H)\leqslant |M_{21}|(n/6 + O(\gamma^{1/3}n)).
$$
It follows from $\gamma\ll 1$ that $|M_{21}|>2g+f-2$, and hence $|M_{21}|\geqslant 2g+f-1$. Similarly, by considering the in-degrees of vertices in $V_3$, we obtain $|M_{13}|\geqslant g-f-1$.

Next we claim that there exists a $D_{1,b,c}$-tiling $\mathcal{H}_1$ such that
$$
|V_1\backslash V(\mathcal{H}_1)|\equiv|V_2\backslash V(\mathcal{H}_1)|\equiv|V_3\backslash V(\mathcal{H}_1)|\pmod h.
$$
If $g=0$, then $|V_1|=|V_2|=|V_3|$, and there is nothing to prove. Thus we may assume that $g\geqslant 1$. 

When $f\geqslant 0$, since
$$
|M_{21}|\geqslant 2g+f-1\geqslant g+f\geqslant R_h(g)+R_h(f),
$$
we may first choose $R_h(g)$ edges from $M_{21}$ and then choose another $R_h(f)$ edges from the remaining edges of $M_{21}$. By Lemma \ref{LEM-edgegood}, we obtain that the former $R_h(g)$ edges have vertex sets that are $\gamma$-transfer for $(V_1,V_3)$, while the latter $R_h(f)$ edges have vertex sets that are $\gamma$-transfer for $(V_2,V_3)$.

When $f<0$, we have
$$
|M_{13}|\geqslant g-f-1\geqslant -f\geqslant R_h(-f).
$$
By Lemma~\ref{LEM-edgegood}, we may choose $R_h(-f)$ edges from $M_{13}$ whose vertex sets are $\gamma$-transfer for $(V_1,V_2)$. Moreover,
$$
|M_{21}|\geqslant 2g+f-1\geqslant g=(g+f)+(-f)\geqslant R_h(g+f)+R_h(-f),
$$
so there exist $R_h(g+f)$ edges in $M_{21}$ that are disjoint from the previously chosen $R_h(-f)$ edges in $M_{13}$. The vertex sets of these edges are $\gamma$-transfer for $(V_1,V_3)$.

By Definition~\ref{DEF:robust}, in either case the edges chosen above can be extended to a constant-sized $D_{1,b,c}$-tiling $\mathcal H_1$ such that
$$
|V_1\setminus V(\mathcal H_1)|\equiv |V_2\setminus V(\mathcal H_1)|\equiv |V_3\setminus V(\mathcal H_1)|\equiv kh/3 \pmod h
$$
for some $k\in\{0,1,2\}$. Moreover, $
(V_1\setminus V(\mathcal H_1),\,V_2\setminus V(\mathcal H_1),\,V_3\setminus V(\mathcal H_1))
$
is still a $\gamma^{1/3}$-superextremal partition of $V(G)\setminus V(\mathcal H_1)$.

Recall that $h$ is divisible by $3$. Hence $kh/3$ is an integer. Let
\begin{equation*}
(x,y,z) = \begin{cases}
(kh/3, 0, 2kh/3), & \mbox{if } b\equiv0\pmod{3}; \\
(kh/3, 0, 0), & \mbox{if } b\equiv1\pmod{3}; \\
(kh/3, 2kh/3, 0), & \mbox{if } b\equiv2\pmod{3}.
\end{cases}
\end{equation*}
Note that $x,y$, and $z$ are small integers. By Lemma \ref{LEM:greedily Dabc}, there is a collection $\mathcal{H}_2$ of disjoint copies of $D_{1,b,c}$ in $G- V(\mathcal{H}_1)$, consisting of $x$ copies of $D_{1,b,c}$ of type-$(1,2,3)$, $y$ copies of $D_{1,b,c}$ of type-$(2,3,1)$, and $z$ copies of $D_{1,b,c}$ of type-$(3,1,2)$, respectively. Next we claim that
$$
|V_i\cap V(\mathcal{H}_2)| \equiv kh/3\pmod h
\quad\text{for each } i\in[3].
$$
Here we only consider the case $b\equiv0\pmod{3}$; the other two cases can be verified similarly. 

Since $3\mid(1+b+c)$ and $b\equiv0\pmod{3}$, it follows that $c\equiv2\pmod3$. Recall that in this case $\mathcal{H}_2$ consists of $kh/3$ copies of $D_{1,b,c}$ of type-$(1,2,3)$ and $2kh/3$ copies of $D_{1,b,c}$ of type-$(3,1,2)$. As $b\equiv0\pmod{3}$, we have $2bk/3\in\mathbb{Z}$, and hence $2bkh/3\equiv0\pmod h$. This implies that $
|V_1\cap V(\mathcal{H}_2)|=kh/3+2bkh/3\equiv kh/3\pmod h$. 
Similarly, $|V_2\cap V(\mathcal{H}_2)|=(b+2c)kh/3\equiv kh/3\pmod{h}$ 
and $|V_3\cap V(\mathcal{H}_2)|=(c+2)kh/3\equiv kh/3\pmod{h}$, 
since $b+2c\equiv c+2\equiv1\pmod3$. 

Therefore, $\mathcal{H}=\mathcal{H}_1\cup\mathcal{H}_2$ is the desired $D_{1,b,c}$-tiling. Let $V_i^\prime=V_i\backslash V(\mathcal{H})$. Clearly, $|V_i^\prime|\equiv0 \pmod h$. Moreover, since $|V(\mathcal{H})|$ is constant, it is not difficult to check that $(V_1^\prime,V_2^\prime,V_3^\prime)$ is a $\gamma^{1/3}$-superextremal partition of $V(G)\backslash V(\mathcal{H})$.
\end{proof}

\end{document}
